% GENERATED FILE. Edit canonical lessons, then run npm run generate:books.
% canonical-source-hash: fnv1a64:8a7abef63e9cea12
% canonical-lessons: GU-C459-komal, GU-C459-kathor, GU-C459-krur, GU-C459-bhalu, GU-C459-pigalvu

\chapter{Gentle and Harsh}
\label{ch:gu-gentle-and-harsh}
\hlchaptermodality{459}

\begin{chapteropening}
\textbf{By the end of this chapter:} \emph{I can describe how gently or harshly someone speaks.}

Complete the last lesson of chapter 459: I can describe how gently or harshly someone speaks.
\end{chapteropening}

\section[\texorpdfstring{\gu{કોમળ} (komaḷ)}{કોમળ (komaḷ)}]{\gu{કોમળ} (komaḷ) — gentle; tender}

\label{lesson:GU-C459-komal}



\begin{quote}
Before the new one: say the Gujarati for to crave, then the Gujarati for to come to mind.
\end{quote}

\subsection*{You'll want to know: \gu{કોમળ}}
\textbf{\gu{કોમળ}} — \emph{komaḷ} — \textquotedblleft{}gentle; tender\textquotedblright{}.

\begin{quote}
\gu{નર્સનો} \gu{અવાજ} \gu{ઘણો} \gu{કોમળ} \gu{હતો}.
\end{quote}

\begin{quote}
\emph{narsano avāj ghaṇo komaḷ hato.}
\end{quote}

\begin{quote}
The nurse's voice was very gentle.
\end{quote}

\begin{cousinweb}
\textbf{\gu{કોમળ}} (\emph{komaḷ}) goes back to Sanskrit \emph{kōmala}, \textquotedblleft{}tender, soft\textquotedblright{}; CDIAL lists the Gujarati word, \textquotedblleft{}soft\textquotedblright{} (CDIAL 3523).
\end{cousinweb}

\begin{grammarlens}[title={What you've built}]
One more word for describing what you lived and felt.
\end{grammarlens}

\subsection*{Your turn}
\begin{itemize}
\raggedright
  \item \emph{Say it:} \emph{komaḷ}
  \item \emph{Say it:} \emph{komaḷ}, once more
  \item \emph{Say it:} \emph{narsano avāj ghaṇo komaḷ hato.} — The nurse's voice was very gentle.
  \item \emph{Recall:} say the Gujarati for to crave, then the Gujarati for to come to mind, then say \emph{komaḷ} again
\end{itemize}

\begin{tcolorbox}[breakable,colback=teal!4,colframe=teal!35!black,title={Before you move on}]
What does \gu{કોમળ} mean? (\textquotedblleft{}Gentle; tender\textquotedblright{}.) Say it once more.
\end{tcolorbox}
\section[\texorpdfstring{\gu{કઠોર} (kaṭhor)}{કઠોર (kaṭhor)}]{\gu{કઠોર} (kaṭhor) — harsh; stern}

\label{lesson:GU-C459-kathor}



\begin{quote}
Before the new one: say the Gujarati for to come to mind, then the Gujarati for tender.
\end{quote}

\subsection*{You'll want to know: \gu{કઠોર}}
\textbf{\gu{કઠોર}} — \emph{kaṭhor} — \textquotedblleft{}harsh; stern\textquotedblright{}.

\begin{quote}
\gu{પિતાએ} \gu{કઠોર} \gu{અવાજે} \gu{જવાબ} \gu{આપ્યો}.
\end{quote}

\begin{quote}
\emph{pitāe kaṭhor avāje javāb āpyo.}
\end{quote}

\begin{quote}
Father answered in a harsh voice.
\end{quote}

\begin{cousinweb}
\textbf{\gu{કઠોર}} (\emph{kaṭhor}) is a tatsama, Sanskrit \emph{kaṭhora}, \textquotedblleft{}hard, solid, stiff\textquotedblright{}, as Monier-Williams records it.
\end{cousinweb}

\begin{grammarlens}[title={What you've built}]
One more word for describing what you lived and felt.
\end{grammarlens}

\subsection*{Your turn}
\begin{itemize}
\raggedright
  \item \emph{Say it:} \emph{kaṭhor}
  \item \emph{Say it:} \emph{kaṭhor}, once more
  \item \emph{Say it:} \emph{pitāe kaṭhor avāje javāb āpyo.} — Father answered in a harsh voice.
  \item \emph{Recall:} say the Gujarati for to come to mind, then the Gujarati for tender, then say \emph{kaṭhor} again
\end{itemize}

\begin{tcolorbox}[breakable,colback=teal!4,colframe=teal!35!black,title={Before you move on}]
What does \gu{કઠોર} mean? (\textquotedblleft{}Harsh; stern\textquotedblright{}.) Say it once more.
\end{tcolorbox}
\section[\texorpdfstring{\gu{ક્રૂર} (krūr)}{ક્રૂર (krūr)}]{\gu{ક્રૂર} (krūr) — cruel}

\label{lesson:GU-C459-krur}



\begin{quote}
Before the new one: say the Gujarati for tender, then the Gujarati for harsh.
\end{quote}

\subsection*{You'll want to know: \gu{ક્રૂર}}
\textbf{\gu{ક્રૂર}} — \emph{krūr} — \textquotedblleft{}cruel\textquotedblright{}.

\begin{quote}
\gu{પ્રાણીઓને} \gu{મારવું} \gu{ક્રૂર} \gu{કામ} \gu{છે}.
\end{quote}

\begin{quote}
\emph{prāṇīone mārvũ krūr kām che.}
\end{quote}

\begin{quote}
Beating animals is a cruel thing to do.
\end{quote}

\begin{cousinweb}
\textbf{\gu{ક્રૂર}} (\emph{krūr}) is a tatsama, Sanskrit \emph{krūra}, \textquotedblleft{}bloody, raw\textquotedblright{} and \textquotedblleft{}cruel, pitiless\textquotedblright{}, as Monier-Williams records it.
\end{cousinweb}

\begin{grammarlens}[title={What you've built}]
One more word for describing what you lived and felt.
\end{grammarlens}

\subsection*{Your turn}
\begin{itemize}
\raggedright
  \item \emph{Say it:} \emph{krūr}
  \item \emph{Say it:} \emph{krūr}, once more
  \item \emph{Say it:} \emph{prāṇīone mārvũ krūr kām che.} — Beating animals is a cruel thing to do.
  \item \emph{Recall:} say the Gujarati for tender, then the Gujarati for harsh, then say \emph{krūr} again
\end{itemize}

\begin{tcolorbox}[breakable,colback=teal!4,colframe=teal!35!black,title={Before you move on}]
What does \gu{ક્રૂર} mean? (\textquotedblleft{}Cruel\textquotedblright{}.) Say it once more.
\end{tcolorbox}
\section[\texorpdfstring{\gu{ભલું} (bhalũ)}{ભલું (bhalũ)}]{\gu{ભલું} (bhalũ) — good-hearted; kind}

\label{lesson:GU-C459-bhalu}



\begin{quote}
Before the new one: say the Gujarati for harsh, then the Gujarati for cruel.
\end{quote}

\subsection*{You'll want to know: \gu{ભલું}}
\textbf{\gu{ભલું}} — \emph{bhalũ} — \textquotedblleft{}good-hearted; kind\textquotedblright{}.

\begin{quote}
\gu{અમારા} \gu{પડોશી} \gu{ભલા} \gu{માણસ} \gu{છે}.
\end{quote}

\begin{quote}
\emph{amārā paḍośī bhalā māṇas che.}
\end{quote}

\begin{quote}
Our neighbour is a kind man.
\end{quote}

\begin{cousinweb}
\textbf{\gu{ભલું}} (\emph{bhalũ}) is inherited from Sanskrit \emph{bhalla}, \textquotedblleft{}auspicious\textquotedblright{}, a rustic form of \emph{bhadra}, \textquotedblleft{}good\textquotedblright{}; CDIAL lists the Gujarati word (CDIAL 9408).
\end{cousinweb}

\begin{grammarlens}[title={What you've built}]
One more word for describing what you lived and felt.
\end{grammarlens}

\subsection*{Your turn}
\begin{itemize}
\raggedright
  \item \emph{Say it:} \emph{bhalũ}
  \item \emph{Say it:} \emph{bhalũ}, once more
  \item \emph{Say it:} \emph{amārā paḍośī bhalā māṇas che.} — Our neighbour is a kind man.
  \item \emph{Recall:} say the Gujarati for harsh, then the Gujarati for cruel, then say \emph{bhalũ} again
\end{itemize}

\begin{tcolorbox}[breakable,colback=teal!4,colframe=teal!35!black,title={Before you move on}]
What does \gu{ભલું} mean? (\textquotedblleft{}Good-hearted; kind\textquotedblright{}.) Say it once more.
\end{tcolorbox}
\section[\texorpdfstring{\gu{પીગળવું} (pīgaḷvũ)}{પીગળવું (pīgaḷvũ)}]{\gu{પીગળવું} (pīgaḷvũ) — to soften with pity; to melt}

\label{lesson:GU-C459-pigalvu}



\begin{quote}
Before the new one: say the Gujarati for cruel, then the Gujarati for good-hearted.
\end{quote}

\subsection*{You'll want to know: \gu{પીગળવું}}
\textbf{\gu{પીગળવું}} — \emph{pīgaḷvũ} — \textquotedblleft{}to soften with pity; to melt\textquotedblright{}.

\begin{quote}
\gu{બાળકનાં} \gu{આંસુ} \gu{જોઈને} \gu{પિતાનું} \gu{હૃદય} \gu{પીગળી} \gu{ગયું}.
\end{quote}

\begin{quote}
\emph{bāḷaknā̃ ā̃su joīne pitānũ hṛday pīgḷī gayũ.}
\end{quote}

\begin{quote}
Seeing the child's tears, the father's heart melted.
\end{quote}

Past (it …): \gu{તે} \gu{બરફ} \gu{પીગળી} \gu{ગયો}, \emph{te baraph pīgḷī gayo}.

\begin{cousinweb}
\textbf{\gu{પીગળવું}} (\emph{pīgaḷvũ}) is inherited from Sanskrit \emph{pragalati}, \textquotedblleft{}drips\textquotedblright{}; CDIAL lists the Gujarati verb, \textquotedblleft{}to melt, be dissolved\textquotedblright{} (CDIAL 8468).
\end{cousinweb}

\begin{grammarlens}[title={What you've built}]
One more word for describing what you lived and felt.
\end{grammarlens}

\subsection*{Your turn}
\begin{itemize}
\raggedright
  \item \emph{Say it:} \emph{pīgaḷvũ}
  \item \emph{Say it:} \emph{pīgaḷvũ}, once more
  \item \emph{Say it:} \emph{bāḷaknā̃ ā̃su joīne pitānũ hṛday pīgḷī gayũ.} — Seeing the child's tears, the father's heart melted.
  \item \emph{Recall:} say the Gujarati for cruel, then the Gujarati for good-hearted, then say \emph{pīgaḷvũ} again
\end{itemize}

\begin{tcolorbox}[breakable,colback=teal!4,colframe=teal!35!black,title={Before you move on}]
What does \gu{પીગળવું} mean? (\textquotedblleft{}To soften with pity; to melt\textquotedblright{}.) Say it once more.
\end{tcolorbox}
